\section{Discussione}
I risultati sperimentali consentono una lettura unificata del sistema come dispositivo di coordinamento passivo incorporato nello spazio. Ciascun esperimento rafforza un collegamento specifico con la letteratura.

\subsection{Collegamento con Burstedde}
L'architettura del modello conferma l'efficacia del paradigma a floor field introdotto da Burstedde et al. \cite{burstedde2001}. I campi statici e dinamici riescono a descrivere in modo compatto, discreto e computazionalmente trasparente il comportamento collettivo dei pedoni in una topologia non banale.

\subsection{Collegamento con Zhang}
L'aumento del tempo medio di uscita al crescere della densità è in linea con il quadro sperimentale dei \textit{fundamental diagrams} in geometrie vincolate \cite{zhang2012}. Il dato interessante è che la rotonda sembra ritardare il collasso rispetto a un'intersezione lineare, ma non ne elimina la possibilità.

\subsection{Collegamento con Helbing}
Il fatto che $\phi$ resti massimo anche per $K_d = 0$ suggerisce una differenza sostanziale rispetto ai casi in cui lane formation emerge soltanto da interazioni locali ripetute \cite{helbing2005}. Nel nostro modello, l'ordine è fortemente geometry-driven: la topologia della rotonda comprime lo spazio delle soluzioni e rende l'ordine collettivo quasi inevitabile.

\subsection{Collegamento con Weng e Yanagisawa}
I risultati sul raggio della rotonda sono coerenti con l'idea, presente in Weng et al., che la rotonda migliori la gestione dei flussi incrociati \cite{weng2006}. Tuttavia, la riduzione monotona del throughput all'aumentare del raggio è in linea soprattutto con la lettura di Yanagisawa et al. \cite{yanagisawa2019}: il central obstacle è utile solo finché non introduce un costo spaziale superiore al beneficio di separazione.

\subsection{L'etologia dei Social Types come limite sistemico}
L'aggiunta dell'approccio a \textit{Psicologia Dinamica} dimostra il limite passivo della rotonda, come visto nei risultati sui Social Types dell'Esperimento 5. Un sistema distribuito fondato sulla repulsione spaziale e sull'allineamento è molto sensibile all'estremizzazione dei vincoli. La tendenza egoistica massimizzata degli Aggressivi, ad esempio, garantisce una percorrenza migliore per il singolo, innescando però parallelamente \textit{clogging} periferici a livello micro-locale e degradando la transizione laminare che la rotonda vorrebbe altrimenti imporre in perfetta stigmergia.

\subsection{Limiti del modello}
Il modello resta un automa cellulare discreto. Questo comporta almeno tre limiti:
\begin{itemize}
\item spazio e tempo sono quantizzati;
\item la repulsione è espressa come costo locale, non come forza continua;
\item la percezione degli agenti non include anisotropie visive o anticipazione avanzata.
\end{itemize}

Di conseguenza, il modello è eccellente per studiare \textit{throughput}, esclusione spaziale, segregazione e topologia del coordinamento, ma meno adatto a descrivere micro-dinamiche cinematiche continue.
